\subsection[Counting Sort]{Counting Sort\cite[p.194~196]{cormen2009introduction} \cite{gfg_countingsort}}
    \subsubsection{Overview}
        \begin{itemize}
            \item \textbf{History and Origin:} In 1954, Harold H. Seward introduced the idea of distribution counting for radix sort\cite[p.~170]{knuth1998sorting}, which later became the basis of the counting sort algorithm.
            \item \textbf{Efficiency:} This algorithm works best when the range of values is smaller than the number of elements. It cannot be used directly for decimal values or non-integer numbers. If the value range is too large, it may consume more memory compared to other in-place algorithms.
            \item \textbf{Application:} Counting sort is often used to optimize radix sort. Its core idea is also applied in bucket sort. It is suitable for sorting data with a limited range of values, such as sorting students by grades or organizing events by time, days, or months.
        \end{itemize}
    \subsubsection{Core Concepts}
        \begin{itemize}
            \item Instead of comparing elements, counting sort counts how many times each value appears in the array.
            \item Then it calculates the sum of previous counts to know the position of each value in the sorted array.
            \item Finally, the elements are placed into their correct positions based on these counts.
        \end{itemize}
    \subsubsection{Step-by-step Explanation}
        \begin{itemize}
            \item Consider an array of n elements, $a[0 \ldots (n-1)]$.
            \item \textbf{Step 1:} Find $k$, the maximum value in array $a$, and create an array $Cnt[0 \ldots k]$.
            \item \textbf{Step 2:} Count the frequency of each value in array $a$ and store it in $Cnt$.
            \item \textbf{Step 3:} Compute the sum of previous counts in $Cnt$ to know the position of each value (prefix sum).
            \item \textbf{Step 4:} Start from $n-1$ and move to $0$, placing each element into its correct position and decreasing the corresponding count in $Cnt$.        \end{itemize}
    \subsubsection{Pseudo Code}
        \begin{itemize}
            \item \textbf{Algorithm:} Counting Sort
            \item \textbf{Input:} An array A of n elements
        \end{itemize}
        \begin{verbatim}
        1. maxv ← 0
        2. for i ← 0 to length(A)-1
            if A[i] > maxv
                    maxv ← A[i]
        3. create Count[0..maxv] and set all to 0
        4. for i ← 0 to length(A)-1
            Count[A[i]] ← Count[A[i]] + 1
        5. for i ← 1 to maxv
            Count[i] ← Count[i] + Count[i-1]
        6. create array B[0..length(A)-1]
        7. for i ← length(A)-1 downto 0
            B[Count[A[i]]-1] ← A[i]
            Count[A[i]] ← Count[A[i]] - 1
        8. for i ← 0 to length(A)-1
            A[i] ← B[i]
        \end{verbatim}
    \subsubsection{Complexity Analysis}
        \begin{itemize}
            \item \textbf{Time Complexity:}
            \begin{itemize}
                \item Best, Average, and Worst case: $O((n + k))$
                    \subitem Where \textit{n} is the number of elements, and \textit{k} is the maximum value in the array.
            \end{itemize}
            \item \textbf{Space Complexity:} $O(n + k)$ 
                \subitem Where \textit{n} is the number of elements, and \textit{k} is the maximum value in the array.
        \end{itemize}
        
    \subsubsection{Variants and Optimizations}
        \begin{itemize}
        \item \textbf{Unstable counting sort\cite{bajpai2014ecounting}:}
            \begin{itemize}
                \item Instead of calculating the prefix sum, we find the correct position directly using the $Cnt$ array. This can make the algorithm run faster but it loses stability.
                \item The time complexity and space complexity are the same as the standard counting sort.
            \end{itemize}
        \end{itemize}
        \begin{itemize}
        \item \textbf{Counting sort with negative numbers\cite{stackoverflow_countingsort_negative}:}
            \begin{itemize}
                \item \textbf{Idea 1:} Find the maximum and minimum values in the array, then create an array with size $(maximum - minimum + 1)$. Instead of $Cnt[a[i]]$, use $Cnt[a[i] - minimum]$. All other steps are the same as the standard counting sort.
                \item \textbf{Idea 2:} Increase the size of the $Cnt$ array to cover the range $[-max, max]$. Instead of $Cnt[a[i]]$, use $Cnt[a[i] + max]$. All other steps remain the same.
                \item \textbf{Idea 3\cite{saiana2018parallelcounting}:} Separate the array into two lists: negative and positive numbers. Convert the negative numbers to positive values (for example by multiplying by $-1$) or apply \textbf{Idea 1} or \textbf{Idea 2}.
            \end{itemize}
        \end{itemize}

